\subsection{FIELDS}

DEFINITION 1. A ring K is called a field if it does not consist only of 0 and every non-zero element of K is invertible.

The set of non-zero elements of the field K, with multiplication, is a group, which is precisely the multiplicative group \(K^{*}\) of the ring K ( \(\S\) 8, no. 1). The opposite ring of a field is a field. A field is called commutative if its multiplication is commutative; such a field is identified with its opposite. A non-commutative field is sometimes called a skew field.

Examples. *(1) We shall define in no. 4 the field of rational numbers; In General Topology the field of real numbers will be defined (General Topology, IV, § 1, no. 3), also the field of complex numbers (General Topology, VIII, § 1, no. 1) and the field of quaternions (General Topology, VIII, § 1, no. 4). These fields are commutative with the exception of the field of quaternions.* (2) The ring Z/2Z is obviously a field.

Let K be a field. Every subring L of K which is a field is called a subfield of K and K is then called an extension field of L; it amounts to the same to say that L is a subring of K and that \(x^{-1} \in L\) for every non-zero element x of L. If \((\mathbf{L}_{i})_{i \in I}\) is a family of subfields of K, then \(\bigcap_{i \in I} L_{i}\) is a subfield of K; for every subset X of K there thus exists a smallest subfield of K containing X; it is said to be generated by X.

PROPOSITION 1. Let K be a field. For every subset X of K, the centralizer (§ 8, no. 5, Example 3) X' of X is a subfield of K.

We know (loc. cit.) that \(\mathbf{X}'\) is a subring of \(\mathbf{K}\) . On the other hand, if \(x \neq 0\) is permutable with \(z \in \mathbf{X}\) , so is \(x^{-1}\) (I, §2, no. 3, Proposition 5) and hence \(\mathbf{X}'\) contains the inverse of every non-zero element of \(\mathbf{X}'\) .

COROLLARY. The centre of a field K is a (commutative) subfield of K.

THEOREM 1. Let A be a ring. The following conditions are equivalent:

\begin{enumerate}
\def\labelenumi{(\alph{enumi})}
\item
  A is a field;
\item
  A is not reduced to 0 and the only left ideals of A are 0 and A.
\end{enumerate}

Suppose that A is a field. Then A is not reduced to 0. Let a be a left ideal of A distinct from 0. There exists a non-zero a belonging to a. For all \(x \in A\) , \(x = (xa^{-1})a \in a\) ; hence a = A.

Suppose that A satisfies condition (b). Let \(x \neq 0\) be in A. We need to prove that x is invertible. The left ideal Ax contains x and hence is not zero, whence Ax = A. There thus exists \(x' \in A\) such that \(x'x = 1\) . Then \(x' \neq 0\) since \(1 \neq 0\) . Applying the above result to \(x'\) , it is seen that there exists \(x'' \in A\) such that \(x''x' = 1\) . Then \(x'' = x'' \cdot 1 = x''x'x = 1 \cdot x = x\) , hence \(xx' = 1\) . Hence \(x'\) is the inverse of x.

Remark. In Theorem 1, left ideals may be replaced by right ideals. In Chapter VIII, § 5, no. 2, we shall study non-zero rings A which have no two-sided ideal distinct from 0 and A; such rings (called quasi-simple) are not necessarily fields * (for example, the ring \(\mathbf{M}_2(\mathbf{Q})\) of square matrices of order 2 with rational coefficients is quasi-simple but is not a field)*.

COROLLARY 1. Let A be a ring and a two-sided ideal of A. For the ring A/a to be a field, it is necessary and sufficient that a be a maximal left ideal of A.

The left ideals of A/a are of the form b/a where b is a left ideal of A containing a (§ 8, no. 8, Corollary to Proposition 5). To say that A/a ≠ 0 means that a ≠ A. Under this hypothesis, A/a is a field if and only if the only left ideals of A containing a are a and A (Theorem 1), whence the corollary.

COROLLARY 2. Let A be a commutative ring which is not 0. There exists a homomorphism of A onto a commutative field.

By Krull's Theorem (§ 8, no. 6, Theorem 1), there exists in A a maximal ideal \(a\) . Then A/a is a field (Corollary 1).

COROLLARY 3. Let a be an integer \(\geqslant 0\) . For the ring \(\mathbf{Z} / a\mathbf{Z}\) to be a field, it is necessary and sufficient that a be prime.

This follows from Corollary 1 and § 8, no. 11.

For p prime the field Z/pZ is denoted by \(F_{p}\) .

THEOREM 2. Let K be a field and A a non-zero ring. If f is a homomorphism of K into A, then the subring \(f(\mathbf{K})\) of A is a field and f defines an isomorphism of K onto \(f(\mathbf{K})\) .

Let \(\mathfrak{a}\) be the kernel of \(f\) . Then \(1 \notin \mathfrak{a}\) for \(f(1) = 1 \neq 0\) in \(\mathbf{A}\) and, as \(\mathfrak{a}\) is a left ideal of \(\mathbf{K}\) , \(\mathfrak{a} = \{0\}\) by Theorem 1. Therefore \(f\) is injective and hence an isomorphism of \(\mathbf{K}\) onto the subring \(f(\mathbf{K})\) of \(\mathbf{A}\) ; the latter ring is therefore a field.

